% Table Style — Booktabs formatting with color highlighting macros
% Usage: Include this file in your paper preamble with % Table Style — Booktabs formatting with color highlighting macros
% Usage: Include this file in your paper preamble with % Table Style — Booktabs formatting with color highlighting macros
% Usage: Include this file in your paper preamble with % Table Style — Booktabs formatting with color highlighting macros
% Usage: Include this file in your paper preamble with \input{table-style}

% Required packages (include in your main .tex preamble if not already present):
% \usepackage{booktabs}
% \usepackage{xcolor}
% \usepackage{colortbl}

% ==============================================================================
% Color Definitions
% ==============================================================================

% Highlight colors for best/second-best results
\definecolor{bestred}{HTML}{FF0000}
\definecolor{secondblue}{HTML}{0000FF}

% Soft row shading (optional)
\definecolor{rowgray}{gray}{0.95}

% ==============================================================================
% Result Highlighting Macros
% ==============================================================================

% \best{value} — Bold red for best result
\newcommand{\best}[1]{\textcolor{bestred}{\textbf{#1}}}

% \second{value} — Underlined blue for second-best result
\newcommand{\second}[1]{\textcolor{secondblue}{\underline{#1}}}

% \gain{value} — Green for improvement/gain values
\definecolor{gaingreen}{HTML}{008000}
\newcommand{\gain}[1]{\textcolor{gaingreen}{#1}}

% ==============================================================================
% Example: Comparison Table
% ==============================================================================
%
% \begin{table}[t]
% \centering
% \caption{Quantitative comparison on [Dataset].
%          \best{Red}: best, \second{blue}: second best.}
% \label{tab:comparison}
% \begin{tabular}{lccc}
% \toprule
% Method & PSNR $\uparrow$ & SSIM $\uparrow$ & LPIPS $\downarrow$ \\
% \midrule
% Baseline A~\cite{a}  & 28.34 & 0.912 & 0.087 \\
% Baseline B~\cite{b}  & \second{29.12} & \second{0.934} & \second{0.065} \\
% \midrule
% Ours                  & \best{31.05} & \best{0.958} & \best{0.042} \\
% \bottomrule
% \end{tabular}
% \end{table}
%
% ==============================================================================
% Example: Ablation Table
% ==============================================================================
%
% \begin{table}[t]
% \centering
% \caption{Ablation study on core components.}
% \label{tab:ablation}
% \begin{tabular}{lccc}
% \toprule
% Configuration & Metric 1 & Metric 2 & Metric 3 \\
% \midrule
% w/o Module A  & 27.8 & 0.89 & 0.102 \\
% w/o Module B  & 29.1 & 0.92 & 0.078 \\
% w/o Module C  & 28.5 & 0.91 & 0.085 \\
% \midrule
% Full model    & \best{31.1} & \best{0.96} & \best{0.042} \\
% \bottomrule
% \end{tabular}
% \end{table}
%
% ==============================================================================
% Guidelines (from paper-writing skill):
% ==============================================================================
% 1. Caption ABOVE the table
% 2. NO vertical lines
% 3. Use \toprule, \midrule, \bottomrule (booktabs)
% 4. Minimize horizontal lines
% 5. Add color to highlight best numbers
% 6. Use arrows (↑/↓) to indicate metric direction
% 7. Single-column tables look better in the right column


% Required packages (include in your main .tex preamble if not already present):
% \usepackage{booktabs}
% \usepackage{xcolor}
% \usepackage{colortbl}

% ==============================================================================
% Color Definitions
% ==============================================================================

% Highlight colors for best/second-best results
\definecolor{bestred}{HTML}{FF0000}
\definecolor{secondblue}{HTML}{0000FF}

% Soft row shading (optional)
\definecolor{rowgray}{gray}{0.95}

% ==============================================================================
% Result Highlighting Macros
% ==============================================================================

% \best{value} — Bold red for best result
\newcommand{\best}[1]{\textcolor{bestred}{\textbf{#1}}}

% \second{value} — Underlined blue for second-best result
\newcommand{\second}[1]{\textcolor{secondblue}{\underline{#1}}}

% \gain{value} — Green for improvement/gain values
\definecolor{gaingreen}{HTML}{008000}
\newcommand{\gain}[1]{\textcolor{gaingreen}{#1}}

% ==============================================================================
% Example: Comparison Table
% ==============================================================================
%
% \begin{table}[t]
% \centering
% \caption{Quantitative comparison on [Dataset].
%          \best{Red}: best, \second{blue}: second best.}
% \label{tab:comparison}
% \begin{tabular}{lccc}
% \toprule
% Method & PSNR $\uparrow$ & SSIM $\uparrow$ & LPIPS $\downarrow$ \\
% \midrule
% Baseline A~\cite{a}  & 28.34 & 0.912 & 0.087 \\
% Baseline B~\cite{b}  & \second{29.12} & \second{0.934} & \second{0.065} \\
% \midrule
% Ours                  & \best{31.05} & \best{0.958} & \best{0.042} \\
% \bottomrule
% \end{tabular}
% \end{table}
%
% ==============================================================================
% Example: Ablation Table
% ==============================================================================
%
% \begin{table}[t]
% \centering
% \caption{Ablation study on core components.}
% \label{tab:ablation}
% \begin{tabular}{lccc}
% \toprule
% Configuration & Metric 1 & Metric 2 & Metric 3 \\
% \midrule
% w/o Module A  & 27.8 & 0.89 & 0.102 \\
% w/o Module B  & 29.1 & 0.92 & 0.078 \\
% w/o Module C  & 28.5 & 0.91 & 0.085 \\
% \midrule
% Full model    & \best{31.1} & \best{0.96} & \best{0.042} \\
% \bottomrule
% \end{tabular}
% \end{table}
%
% ==============================================================================
% Guidelines (from paper-writing skill):
% ==============================================================================
% 1. Caption ABOVE the table
% 2. NO vertical lines
% 3. Use \toprule, \midrule, \bottomrule (booktabs)
% 4. Minimize horizontal lines
% 5. Add color to highlight best numbers
% 6. Use arrows (↑/↓) to indicate metric direction
% 7. Single-column tables look better in the right column


% Required packages (include in your main .tex preamble if not already present):
% \usepackage{booktabs}
% \usepackage{xcolor}
% \usepackage{colortbl}

% ==============================================================================
% Color Definitions
% ==============================================================================

% Highlight colors for best/second-best results
\definecolor{bestred}{HTML}{FF0000}
\definecolor{secondblue}{HTML}{0000FF}

% Soft row shading (optional)
\definecolor{rowgray}{gray}{0.95}

% ==============================================================================
% Result Highlighting Macros
% ==============================================================================

% \best{value} — Bold red for best result
\newcommand{\best}[1]{\textcolor{bestred}{\textbf{#1}}}

% \second{value} — Underlined blue for second-best result
\newcommand{\second}[1]{\textcolor{secondblue}{\underline{#1}}}

% \gain{value} — Green for improvement/gain values
\definecolor{gaingreen}{HTML}{008000}
\newcommand{\gain}[1]{\textcolor{gaingreen}{#1}}

% ==============================================================================
% Example: Comparison Table
% ==============================================================================
%
% \begin{table}[t]
% \centering
% \caption{Quantitative comparison on [Dataset].
%          \best{Red}: best, \second{blue}: second best.}
% \label{tab:comparison}
% \begin{tabular}{lccc}
% \toprule
% Method & PSNR $\uparrow$ & SSIM $\uparrow$ & LPIPS $\downarrow$ \\
% \midrule
% Baseline A~\cite{a}  & 28.34 & 0.912 & 0.087 \\
% Baseline B~\cite{b}  & \second{29.12} & \second{0.934} & \second{0.065} \\
% \midrule
% Ours                  & \best{31.05} & \best{0.958} & \best{0.042} \\
% \bottomrule
% \end{tabular}
% \end{table}
%
% ==============================================================================
% Example: Ablation Table
% ==============================================================================
%
% \begin{table}[t]
% \centering
% \caption{Ablation study on core components.}
% \label{tab:ablation}
% \begin{tabular}{lccc}
% \toprule
% Configuration & Metric 1 & Metric 2 & Metric 3 \\
% \midrule
% w/o Module A  & 27.8 & 0.89 & 0.102 \\
% w/o Module B  & 29.1 & 0.92 & 0.078 \\
% w/o Module C  & 28.5 & 0.91 & 0.085 \\
% \midrule
% Full model    & \best{31.1} & \best{0.96} & \best{0.042} \\
% \bottomrule
% \end{tabular}
% \end{table}
%
% ==============================================================================
% Guidelines (from paper-writing skill):
% ==============================================================================
% 1. Caption ABOVE the table
% 2. NO vertical lines
% 3. Use \toprule, \midrule, \bottomrule (booktabs)
% 4. Minimize horizontal lines
% 5. Add color to highlight best numbers
% 6. Use arrows (↑/↓) to indicate metric direction
% 7. Single-column tables look better in the right column


% Required packages (include in your main .tex preamble if not already present):
% \usepackage{booktabs}
% \usepackage{xcolor}
% \usepackage{colortbl}

% ==============================================================================
% Color Definitions
% ==============================================================================

% Highlight colors for best/second-best results
\definecolor{bestred}{HTML}{FF0000}
\definecolor{secondblue}{HTML}{0000FF}

% Soft row shading (optional)
\definecolor{rowgray}{gray}{0.95}

% ==============================================================================
% Result Highlighting Macros
% ==============================================================================

% \best{value} — Bold red for best result
\newcommand{\best}[1]{\textcolor{bestred}{\textbf{#1}}}

% \second{value} — Underlined blue for second-best result
\newcommand{\second}[1]{\textcolor{secondblue}{\underline{#1}}}

% \gain{value} — Green for improvement/gain values
\definecolor{gaingreen}{HTML}{008000}
\newcommand{\gain}[1]{\textcolor{gaingreen}{#1}}

% ==============================================================================
% Example: Comparison Table
% ==============================================================================
%
% \begin{table}[t]
% \centering
% \caption{Quantitative comparison on [Dataset].
%          \best{Red}: best, \second{blue}: second best.}
% \label{tab:comparison}
% \begin{tabular}{lccc}
% \toprule
% Method & PSNR $\uparrow$ & SSIM $\uparrow$ & LPIPS $\downarrow$ \\
% \midrule
% Baseline A~\cite{a}  & 28.34 & 0.912 & 0.087 \\
% Baseline B~\cite{b}  & \second{29.12} & \second{0.934} & \second{0.065} \\
% \midrule
% Ours                  & \best{31.05} & \best{0.958} & \best{0.042} \\
% \bottomrule
% \end{tabular}
% \end{table}
%
% ==============================================================================
% Example: Ablation Table
% ==============================================================================
%
% \begin{table}[t]
% \centering
% \caption{Ablation study on core components.}
% \label{tab:ablation}
% \begin{tabular}{lccc}
% \toprule
% Configuration & Metric 1 & Metric 2 & Metric 3 \\
% \midrule
% w/o Module A  & 27.8 & 0.89 & 0.102 \\
% w/o Module B  & 29.1 & 0.92 & 0.078 \\
% w/o Module C  & 28.5 & 0.91 & 0.085 \\
% \midrule
% Full model    & \best{31.1} & \best{0.96} & \best{0.042} \\
% \bottomrule
% \end{tabular}
% \end{table}
%
% ==============================================================================
% Guidelines (from paper-writing skill):
% ==============================================================================
% 1. Caption ABOVE the table
% 2. NO vertical lines
% 3. Use \toprule, \midrule, \bottomrule (booktabs)
% 4. Minimize horizontal lines
% 5. Add color to highlight best numbers
% 6. Use arrows (↑/↓) to indicate metric direction
% 7. Single-column tables look better in the right column
